\section{Introduction}\label{sec:overview}

\subsection{The general framework}\label{subsec:framework}

\nick{not accurate subsection title: it is basically a summary of notations, not the proof framework}

An adversary picks a finite multigraph \(\Gin=(\Vset,\Eset)\); parallel edges
and self-loops are allowed. Write
\[
  \nvert:=|\Vset|,
  \qquad
  \medge:=|\Eset|.
\]
Each vertex \(v\) carries a Boolean variable \(\bvar{v}\), and each edge \(uv\)
carries one of the four clauses
\[
  \bvar{u}\vee\bvar{v},
  \qquad
  \bvar{u}\vee\neg\bvar{v},
  \qquad
  \neg\bvar{u}\vee\bvar{v},
  \qquad
  \neg\bvar{u}\vee\neg\bvar{v} .
\]
In other words, each end of the edge gets a sign, which decides whether its
variable appears negated. The two signs of an edge are independent and uniform,
and different edges are independent. A self-loop at \(v\) carries a clause on
\(\bvar{v}\) alone, again with two independent uniform signs. The conjunction
of all \(\medge\) clauses is a random 2-CNF formula, which we denote by
\(\rform{\Gin}\).

Thus the graph is chosen by the adversary, and only the signs are random; this
is what \emph{semi-random} means here. A bound on \(\Prx{\rform{\Gin}\SAT}\)
must therefore hold for every graph that the adversary can pick, and not only
for a typical one.

For a finite graph \(H\), let \(\ncomp{H}\) be its number of connected
components, where each isolated vertex counts as one. Define
\[
  \cyc{H}:=|E(H)|-|V(H)|+\ncomp{H}
\]
to be the total cycle rank of \(H\), that is, the number of independent cycles,
and
\[
  \exc{H}:=|E(H)|-\frac32|V(H)|
\]
to be the excess of \(H\). Positive excess is the same as edge density above
\(3/2\). We write \(\dg{H}{v}\) for the degree of a vertex \(v\) in \(H\); a
self-loop at \(v\) adds \(2\) to it.

For \(v\in V(\Gin)\) and an integer \(\rad\ge0\), let \(\ball{\rad}{v}\) be the
subgraph induced by all vertices \(u\) with graph distance
\(\dist(u,v)\le\rad\).
If the maximum degree of \(\Gin\) is at most \(\degb\), the layer at distance
\(j\ge1\) from \(v\) has at most \(\degb(\degb-1)^{j-1}\) vertices, so
\[
  |V(\ball{\rad}{v})|\le \ballsz{\rad},
  \qquad
  \ballsz{\rad}:=1+\degb\sum_{j=0}^{\rad-1}(\degb-1)^j .
\]
We use \(\ballsz{\rad}\) throughout as the standard bound on the size of a
radius-\(\rad\) ball.

\begin{definition}[Locally tree-like]\label{def:loctree}
Let \(\rad\ge1\) be an integer and \(\nontree>0\). A finite graph \(\Gin\) is
\emph{\loctree{\rad}{\nontree}} if at most \(\nontree|V(\Gin)|\) vertices \(v\)
have a radius-\(\rad\) neighborhood \(\ball{\rad}{v}\) that is not a tree.
\end{definition}

Finally, for a fixed 2-CNF formula \(\fml\), let \(\Sol(\fml)\) be the set of
its satisfying assignments. We always regard \(\Sol(\fml)\) as a graph, the
\emph{solution graph} of \(\fml\), in which two assignments are adjacent when
they differ in exactly one variable. Its number of connected components,
\(\solcomp{\fml}\), is the number of solution clusters of \(\fml\). In this way
\(\ncomp{\cdot}\) always counts the components of a graph, and
\(\solcomp{\fml}=0\) when \(\fml\) is unsatisfiable.

\subsection{Random 2-SAT}

\nick{The section starts as ``relation to random 2-SAT'', then states some of our results and does a warm up lemma (first-moment bound). 3 different things, unclear how are they related}

In the classical random 2-SAT model the graph is random as well: a formula on
\(\nvert\) variables is built from \(\medge\) clauses drawn independently and
uniformly. This model has a sharp threshold at clause density \(1\): the
formula is satisfiable with high probability when \(\medge/\nvert<1\), and
unsatisfiable with high probability when \(\medge/\nvert>1\). The threshold was
found by Chv\'atal--Reed~\cite{ChvatalReed1992} and Goerdt~\cite{Goerdt1996},
and the width of the transition window was determined by Bollob\'as, Borgs,
Chayes, Kim and Wilson~\cite{BBCKW2001}.

Since a clause involves two variables, a 2-CNF formula can be viewed as a
multigraph with a sign on each edge \nick{sign on an edge? you probably meant ``randomly chosen signs of edge variables''}: the vertices are the variables, and the
edges are the clauses. The ratio \(\medge/\nvert\) is then the edge density of
this multigraph. In the classical model the multigraph itself is uniformly
random, so the threshold says that a uniformly random multigraph of density
above \(1\) already makes the formula unsatisfiable.

\nick{the next text up to statement of (Main theorem) seems to be out of scope for this subsection. It makes sense to group it with preliminary results on prob. is greater than 0 if density is higher than $3/2$.}
This paper asks what happens when an adversary chooses the multigraph and only
the signs are random. We show that the threshold then moves from density \(1\)
to density \(3/2\). Up to density \(3/2\), the adversary can choose a graph for
which the formula is satisfiable under every sign pattern. Above \(3/2\), no
choice of graph helps: for every graph of density at least \(3/2+\excpar\), the
formula is unsatisfiable with high probability.

The density \(3/2\) cannot be lowered. Two variables have four assignments, and
a clause on them rules out exactly one of these. Hence three parallel edges
between two vertices always leave at least one satisfying assignment, whatever
the signs are. This graph has two vertices and three edges, so its density is
\(3/2\), and the same is true for any number of disjoint copies of it. Thus
there are graphs with arbitrarily many vertices for which
\[
  \frac{\medge}{\nvert}=\frac32,
  \qquad
  \Prx{\rform{\Gin}\SAT}=1 .
\]
Our main result says that the adversary cannot do better than this.

\begin{theorem}[Main theorem]\label{thm:overview-main}\label{thm:main}
Fix \(\excpar>0\). Then every finite multigraph \(\Gin\) with \(\nvert\)
vertices and
\[
  \exc{\Gin}\ge\excpar\nvert
\]
satisfies
\[
  \Prx{\rform{\Gin}\SAT}\le e^{-\Omega(\nvert)},
\]
where the implied constant depends only on \(\excpar\). In particular the
satisfiability probability tends to \(0\) as \(\nvert\to\infty\).
\end{theorem}

A first attempt is to count satisfying assignments. This gives a bound of the
same form as Theorem~\ref{thm:main} with a very short proof, but only when the
density is at least \(2.41\).

\begin{lemma}[First-moment bound]\label{lem:overview-first-moment}
Every finite multigraph \(\Gin\) with \(\nvert\) vertices and
\(\medge\ge2.41\,\nvert\) edges satisfies
\[
  \Prx{\rform{\Gin}\SAT}\le e^{-\Omega(\nvert)} .
\]
\end{lemma}

\begin{proof}
Fix one assignment. For each edge, exactly one of the four sign choices makes
its clause false under this assignment, so the clause is true with probability
\(3/4\). Since the signs of different edges are independent, the assignment
satisfies \(\rform{\Gin}\) with probability \((3/4)^{\medge}\), and linearity of
expectation over the \(2^{\nvert}\) assignments gives
\(\E\,|\Sol(\rform{\Gin})|=2^{\nvert}(3/4)^{\medge}\). If \(\rform{\Gin}\) is
satisfiable then \(|\Sol(\rform{\Gin})|\ge1\), hence
\[
  \Prx{\rform{\Gin}\SAT}\le\E\,|\Sol(\rform{\Gin})|
  =\exp\!\left[\nvert\log2-\medge\log\tfrac43\right].
\]
Since
\[
  2.41>\frac{\log2}{\log(4/3)}\approx2.409 ,
\]
the exponent is at most
\(-\bigl(2.41\log\tfrac43-\log2\bigr)\nvert\) whenever
\(\medge\ge2.41\,\nvert\), and \(2.41\log\tfrac43-\log2>0\).
\end{proof}

\subsection{Overview of the proof}\label{subsec:overview-proof}

For densities between \(3/2\) and \(2.41\) the first-moment bound fails, and we
need a different argument. The proof has three parts. Here we explain the idea
of each part and how the parts fit together; the next three subsections give
more details.

\medskip
\noindent\textbf{Part 1: reduction to difficult graphs}
(Section~\ref{sec:graph-analysis}). Some graphs are easy. Suppose that
\(\Gin\) contains a large number (linear in $n$) of disjoint subgraphs of constant size, and that
each of them has at least one sign pattern that makes its clauses
unsatisfiable. Then each subgraph gives a contradiction with constant probability. 
Since these events are independent, with probability
\(1-e^{-\Omega(\nvert)}\) at least one of them happens, and then the formula is
unsatisfiable. In Section~\ref{sec:graph-analysis} we show that every graph
either has this structure or can be turned into a \emph{difficult graph}, and
that this change does not decrease the SAT probability. Specifically, the difficult
graph has no loops; its minimum and maximum degree are, respectively, \(\ge 2\) and a constant
\(\le \degb\); its excess is linear in $n$; and the graph is locally tree-like for most of its vertices, i.e., the small neighborhoods (a ball of constant radius) for almost all of its vertices are trees.

\medskip
\noindent\textbf{Part 2: from satisfiability to a random experiment on the
graph} (Section~\ref{sec:solution-components}). For difficult graphs we cannot
rely on the local argument, because almost all small neighborhoods are trees,
and a 2-CNF whose graph is a forest is satisfiable under every sign pattern.
We therefore
need a global argument, which makes this section the most interesting part of the proof. The central idea of this section is to count the number of solution clusters instead of
the number of solutions. This number is at least \(1\) whenever the formula is
satisfiable, so its expectation is an upper bound on
\(\Prx{\rform{\Gin}\SAT}\). The main identity,
Theorem~\ref{thm:overview-main-identity}, computes this expectation exactly in
terms of the following random edge-deletion experiment on \(\Gin\). We go
through the edges of \(\Gin\) in a fixed order and build a forest. 
% If the two endpoints of
% an edge are not yet connected in the forest, we flip a fair coin and add the
% edge to the forest on heads. 
If the two endpoints of
an edge are not yet connected in the forest, we add the
edge to the forest with probability $0.5$ (and remove the edge with the remaining $0.5$ probability). 
If the pair of vertices are already connected, adding the edge
would form a cycle, so we mark this edge and deffer its coin toss.
% deffer the coin toss until  
% do not flip its coin and call the edge deferred.
Let
\(\Tdef\) be the number of deferred edges, and let \(\Kdef\) be the number of
connected components of the final forest that contain a deferred edge (the precise
definitions are given in Section~\ref{subsec:overview-main-identity}). The main
identity says that the expected number of solution clusters \(\Ex{\solcomp{\fml}}\) equals
\(\Ex{2^{\Kdef}(3/4)^{\Tdef}}\). On the left-hand side the expectation is over
the random signs of \(\fml\), while on the right-hand side it is over the random coins of the
experiment with only one coin toss per each edge (instead of two coin tosses per clause in \(\fml\)). The main identity gives
a purely graph-theoretic view of the problem: to bound the satisfiability
probability, we only need to study \(\Kdef\) and \(\Tdef\) on the graph
\(\Gin\).

\medskip
\noindent\textbf{Part 3: the bounds on \(\Kdef\) and \(\Tdef\)}
(Section~\ref{sec:difficult-graphs}). For a difficult graph we show that
\(\Kdef\) is small for every outcome of the experiment, while \(\Tdef\) is
large with high probability. A forest component that contains a deferred edge
also contains a cycle of \(\Gin\), formed by this edge and the forest path
between its endpoints, and different components give disjoint cycles. A locally
tree-like graph cannot contain many disjoint cycles, because it has few short
ones and there is no room for many long ones, so \(\Kdef\) is small. For
\(\Tdef\), suppose that we also flip the coins of the deferred edges. Then
every edge is kept with probability \(1/2\), and since the edges kept before
form a forest, every cycle of the kept subgraph uses a deferred edge. This is
where the density \(3/2\) comes in. Density above \(3/2\) means average degree
above \(3\), so a path that enters a vertex can, on average, leave it along
more than \(2\) other edges. If each edge is kept with probability \(1/2\),
then on average more than one of these ways survives, so long paths survive
and the kept subgraph still has linear cycle rank. This makes
\(\Tdef\) large. As a result,
the factor \((3/4)^{\Tdef}\) beats the factor \(2^{\Kdef}\), and the
expectation is exponentially small.

\medskip
Combining the three parts proves Theorem~\ref{thm:main}.

\subsection{The graph-theoretic analysis}\label{subsec:overview-graph}

In this subsection \(\degb<\infty\) is a degree bound and \(\excd>0\) is an
excess density. Both are constants that do not depend on \(\nvert\); their
final values are fixed in Lemma~\ref{lem:overview-bounded-core} and depend only
on \(\excpar\).

\medskip
\noindent\textbf{\Hyp.} \(\Gin\) has no self-loops, every vertex satisfies
\(2\le\dg{\Gin}{v}\le\degb\), and \(\exc{\Gin}\ge\excd\nvert\).

\medskip
We call a graph \emph{difficult} if it satisfies \Hyp{} and is
\loctree{\rad}{\nontree}. This subsection deals with all other graphs: for each
of them we either prove the conclusion directly or reduce it to a difficult
graph. The reduction has four steps.

\medskip
\noindent\textbf{Step 1: bounded degree.} If \(\Gin\) has at least
\(2.41\nvert\) edges, the first-moment bound of
Lemma~\ref{lem:overview-first-moment} already gives the conclusion, so we may
assume that \(\Gin\) has fewer edges.

\begin{lemma}[Bounded-degree reduction]\label{lem:overview-bounded-core}
Fix \(\excpar>0\). There are constants \(\degb<\infty\) and \(\excd>0\)
depending only on \(\excpar\) such that for every graph \(\Gin\) with
\(\nvert\) vertices, fewer than \(2.41\nvert\) edges and
\(\exc{\Gin}\ge\excpar\nvert\), one can construct a finite multigraph \(H\) with
\[
  |V(H)|\ge\frac{\excpar}{\degb}\,\nvert,
  \qquad
  2\le\dg{H}{v}\le\degb,
  \qquad
  \exc{H}\ge2\excd|V(H)|,
\]
and \(\Prx{\rform{\Gin}\SAT}\le\Prx{\rform{H}\SAT}\).
\end{lemma}

First we may delete vertices of degree \(0\) or \(1\), since an isolated
variable does not occur in the formula, and a variable that
occurs in only one clause can always be set to satisfy it. Both deletions
increase the excess, so they cost nothing. Next, each vertex of large degree is
replaced by a path of copies of bounded degree, and the copies are forced to
take the same value. Every split adds one vertex and one edge and therefore
lowers the excess by \(1/2\); if the target degree is large enough, the total
loss is less than half of the initial excess.

The difficult point is that two clauses are needed to force two neighboring
copies to take the same value, while the new graph has only one edge between
them. The proof in Section~\ref{subsec:bounded-core} solves this by first
flipping the variable of each copy at random and then keeping one of the two
clauses at random. This produces exactly the right distribution of signs, and
removing a clause can only make the formula easier to satisfy.

\medskip
\noindent\textbf{Step 2: components of nonpositive excess.}

\begin{lemma}[Positive-excess component reduction]\label{lem:overview-positive-components}
Let \(\Gin\) be a finite multigraph with maximum degree at most \(\degb\).
Delete every connected component \(C\) with \(\exc{C}\le0\), and call the
result \(H\). Then every vertex of \(H\) has the same degree in \(H\) as in
\(\Gin\),
\[
  \exc{H}\ge\exc{\Gin},
  \qquad
  |V(H)|\ge\frac{2}{\degb}\,\exc{\Gin},
\]
every component of \(H\) without a self-loop has cycle rank at least \(3\),
every component with a self-loop has cycle rank at least \(2\), and
\(\Prx{\rform{\Gin}\SAT}\le\Prx{\rform{H}\SAT}\).
\end{lemma}

Since the excess of a graph is the sum of the excesses of its components,
deleting components of nonpositive excess can only increase the excess, and the
degree bound ensures that the remaining graph is not too small. The main gain,
however, is the statement about cycle ranks. For a connected graph \(C\), the
condition \(\exc{C}>0\) is the same as \(\cyc{C}>1+|V(C)|/2\). If \(C\) has no
self-loop, then \(|V(C)|\ge2\), hence \(\cyc{C}>2\), and since the cycle rank is
an integer, \(\cyc{C}\ge3\). If \(C\) has a self-loop, then \(|V(C)|=1\) is
possible, and we only get \(\cyc{C}\ge2\). Steps 3 and 4 need exactly these
two bounds.

\medskip
\noindent\textbf{Step 3: self-loops.} The main identity of
Section~\ref{subsec:overview-main-identity} holds only for graphs without
self-loops, so we must get rid of them. Deleting a self-loop deletes a clause,
which can only make the formula easier to satisfy, but it lowers the excess by
\(1\). This is why Steps 1 and 2 keep the excess density \(2\excd\): we can
afford to lose \(\excd\nvert\) of the excess. If there are at most
\(\excd\nvert\) self-loops, we delete them all, and the excess density goes down
from \(2\excd\) to \(\excd\). This may create vertices of degree \(0\) or \(1\)
and components of nonpositive excess, so we delete them again as in Steps 1
and 2.

If there are more than \(\excd\nvert\) self-loops, we cannot delete them all,
but in this case they help us. A self-loop at \(v\) carries one of the clauses
\(\bvar{v}\), \(\neg\bvar{v}\) and \(\bvar{v}\vee\neg\bvar{v}\), so with
probability \(1/2\) it fixes the value of \(\bvar{v}\), and many such strong
local constraints lead to many local contradictions. More precisely, we show
that linearly many balls allow a contradiction. For this we need a simple test
that tells us when a connected graph has a sign pattern that makes its formula
unsatisfiable. The following lemma gives such a test in terms of the cycle
rank, and it will be used again in Step 4.

\begin{lemma}[Unsatisfiable support]\label{lem:overview-unsat-support}
Let \(\Gin\) be a finite connected multigraph. Then \(\Gin\) supports an
unsatisfiable 2-CNF, that is, there is a sign pattern on its edges for which
the resulting formula is unsatisfiable, in either of the following cases:
\begin{enumerate}[label=(\arabic*),leftmargin=3em]
\item \(\Gin\) has no self-loops and \(\cyc{\Gin}\ge3\);
\item \(\Gin\) has a self-loop and \(\cyc{\Gin}\ge2\).
\end{enumerate}
\end{lemma}

For simple graphs the first case is due to Karve--Hirani~\cite[Corollary
19]{KarveHirani2019}; parallel edges and self-loops need separate arguments,
which we give in Section~\ref{subsec:unsat-support}. The idea is that, with
suitable signs, cycles let a formula carry a forced value around the graph and
come back to a contradiction. A self-loop with suitable signs fixes the value
of its variable by itself, so one more independent cycle is enough. Note that
the lemma only gives sufficient conditions, and without self-loops cycle rank
\(2\) is not always enough: three parallel edges between two vertices have
cycle rank \(2\), but they are satisfiable under every sign pattern.

The next two lemmas, one for Step 3 and one for Step 4, are proved by the same
idea. We start from linearly many seeds: the vertices with a self-loop in
Lemma~\ref{lem:overview-loop-counting}, and disjoint cycles in
Lemma~\ref{lem:overview-finite-counting}. If several seeds lie close together,
a ball around them has enough cycle rank to allow a contradiction. If a seed is
alone but its whole component is small, then a ball contains this component,
and by Step 2 the component has enough cycle rank. In the remaining case the
seed is alone and its component reaches far away, so the seed has a large
neighborhood without other seeds. These neighborhoods are disjoint, so there
are only few such seeds. The radius \(\radp\) of the balls is chosen inside
each lemma; in Lemma~\ref{lem:overview-finite-counting} it is larger than
\(\rad\).

\begin{lemma}[Self-loop reduction]\label{lem:overview-loop-counting}
Fix \(\degb<\infty\) and \(\excd>0\). There are \(\radp<\infty\) and
\(\pack>0\) depending only on them with the following property. Let \(\Gin\) be
an \(\nvert\)-vertex multigraph with maximum degree at most \(\degb\) in which
every component containing a self-loop has cycle rank at least \(2\). If \(\Gin\)
has more than \(\excd\nvert\) self-loops, then at least \(\pack\nvert\) vertices
\(v\) have a ball \(\ball{\radp}{v}\) that contains a self-loop and has cycle
rank at least \(2\).
\end{lemma}

By case~(2) of Lemma~\ref{lem:overview-unsat-support}, each of these balls
allows a contradiction. Here two self-loops close together are already enough,
since a self-loop needs only one more independent cycle.

\medskip
\noindent\textbf{Step 4: local tree structure.} After Step 3 the graph has no
self-loops, and since Step 2 was applied again, every component has cycle rank
at least \(3\). If the graph is
locally tree-like, it is difficult, the reduction is finished, and the graph is
treated in the next two subsections. Otherwise the next lemma gives linearly many balls
without self-loops and with cycle rank at least \(3\), and by case~(1) of
Lemma~\ref{lem:overview-unsat-support} each of them allows a contradiction.

\begin{lemma}[Local-tree reduction]\label{lem:overview-finite-counting}
Fix \(\degb<\infty\), an integer \(\rad\ge1\) and \(\nontree>0\). There are
\(\radp<\infty\) and \(\pack>0\) depending only on them with the following
property. Let \(\Gin\) be an \(\nvert\)-vertex graph without self-loops, with
maximum degree at most \(\degb\), in which every component has cycle rank at
least \(3\). If \(\Gin\) is not \loctree{\rad}{\nontree}, then at least
\(\pack\nvert\) vertices \(v\) have a ball \(\ball{\radp}{v}\) with cycle rank
at least \(3\).
\end{lemma}

If the graph is not locally tree-like, then linearly many vertices see a cycle
nearby, and a greedy choice among them gives linearly many disjoint cycles,
which are the seeds. Here three cycles must lie close together, because two
cycles give only cycle rank \(2\), which is not enough without self-loops. So
a seed with at most one other seed nearby is handled by the component argument.

The vertices given by Lemma~\ref{lem:overview-loop-counting} or
Lemma~\ref{lem:overview-finite-counting} are exactly what the following
statement needs.

\begin{corollary}[Packing local contradictions]\label{cor:overview-packing}
Fix \(\degb<\infty\), an integer \(\radp\ge1\) and \(\pack>0\). Suppose \(\Gin\)
has \(\nvert\) vertices and maximum degree at most \(\degb\), and that at least
\(\pack\nvert\) vertices \(v\) satisfy one of the following:
\begin{enumerate}[label=(\arabic*),leftmargin=3em]
\item \(\ball{\radp}{v}\) has no self-loop and \(\cyc{\ball{\radp}{v}}\ge3\);
\item \(\ball{\radp}{v}\) has a self-loop and \(\cyc{\ball{\radp}{v}}\ge2\).
\end{enumerate}
Then, with the implied constant depending only on \(\degb,\radp,\pack\),
\[
  \Prx{\rform{\Gin}\SAT}\le e^{-\Omega(\nvert)} .
\]
\end{corollary}

To prove the corollary, we choose the centers greedily so that the balls are
disjoint. Each choice removes only a bounded number of candidates, so we get
linearly many disjoint balls. By Lemma~\ref{lem:overview-unsat-support}, each
of these balls has a sign pattern for which its clauses are unsatisfiable, and
since a ball has a bounded number of edges, its random signs equal this pattern
with probability bounded below by a positive constant. Disjoint balls have
disjoint edges, so these events are independent, and the probability that none
of them happens is exponentially small.

\medskip
\noindent\textbf{Combining the steps.} Let \(\Gin\) satisfy
\(\exc{\Gin}\ge\excpar\nvert\). Steps 1 and 2 turn it into a graph with degrees
in \([2,\degb]\) and excess density at least \(2\excd\), in which every
component has cycle rank at least \(3\), or at least \(2\) if it has a
self-loop; neither step decreases the satisfiability probability. Step 3 either
finishes the proof by packing or removes the self-loops, after which the excess
density is at least \(\excd\). Step 4 either finishes the proof by packing or
gives a difficult graph, with the radius and fraction fixed in
Section~\ref{subsec:overview-three-conditions}. Since every step keeps a constant fraction of the vertices, all the
exponential bounds remain exponential in the original \(\nvert\). The details
are in Section~\ref{sec:graph-analysis} and Section~\ref{subsec:main-proof}.

In the end, therefore, we are left with a difficult graph. For such a graph we
cannot rely on the packing argument,
because a ball that is a tree allows no contradiction and almost all balls of
the fixed radius are trees. The next two subsections deal with this case.

\subsection{The expected number of solution clusters}\label{subsec:overview-main-identity}

We first explain why the first-moment bound is weak. It counts satisfying
assignments, but satisfying assignments come in clusters. Suppose that an
assignment satisfies \(\fml\), and that every clause containing some variable
is also satisfied by its other literal. Then flipping this variable gives
another satisfying assignment, and if several such variables share no clause,
they can be flipped in every combination, which gives a whole subcube of
solutions. As a result, the number of satisfying assignments can be
exponentially larger than the number of really different solutions.

Counting clusters avoids this problem. Since \(\solcomp{\fml}\ge1\) whenever
\(\fml\) is satisfiable, we have
\[
  \Prx{\rform{\Gin}\SAT}\le\E\,\solcomp{\rform{\Gin}},
\]
and \(\solcomp{\fml}\) does not depend on the size of the clusters. The
difficulty is that \(\solcomp{\fml}\) depends on the connectivity of the whole
solution graph, so there is no direct way to compute its expectation. The
identity below solves this problem.

Since the graph given by the reduction has no self-loops, we only need the
identity in this case, and throughout this subsection we assume that \(\Gin\)
has no self-loops.

Run the following random edge-deletion experiment on \(\Gin\), independently of
the signs. Fix an order of the edges, start with the empty graph \(\Ffor\), and
look at the edges one at a time. If the endpoints of the current edge are not
yet connected in \(\Ffor\), flip a fair coin right away and add the edge to
\(\Ffor\) if the coin says keep. If the endpoints are already connected, do not
flip the coin yet, but defer the edge to the end. In this way \(\Ffor\) stays a
forest during the whole first stage. Write
\[
  \Tdef:=\text{number of deferred edges},
  \qquad
  \Kdef:=\text{number of forest components containing a deferred edge}.
\]

\begin{theorem}[Main identity]\label{thm:overview-main-identity}
For every finite self-loop-free multigraph \(\Gin\),
\[
  \E_{\Phi}\,\solcomp{\rform{\Gin}}
  =
  \Ex{2^{\Kdef}\left(\frac34\right)^{\Tdef}} .
\]
\end{theorem}

\begin{corollary}\label{cor:overview-sat-bound}
For every finite self-loop-free multigraph \(\Gin\),
\[
  \Prx{\rform{\Gin}\SAT}
  \le
  \Ex{2^{\Kdef}\left(\frac34\right)^{\Tdef}} .
\]
\end{corollary}

\begin{proof}
If \(\rform{\Gin}\) is satisfiable, then \(\solcomp{\rform{\Gin}}\ge1\).
Hence \(\Prx{\rform{\Gin}\SAT}\le\E_{\Phi}\,\solcomp{\rform{\Gin}}\), and
Theorem~\ref{thm:overview-main-identity} gives the claim.
\end{proof}

Theorem~\ref{thm:overview-main-identity} is proved in
Section~\ref{sec:direct-identity}. The proof rests on three observations, which
we now describe one by one.

The first observation is that the cluster count of a 2-CNF is an Euler
characteristic. View the solution set \(\Sol(\fml)\) inside the Boolean cube
\(\{0,1\}^{\Vset}\), and record not only the satisfying assignments but also
every subcube all of whose assignments satisfy \(\fml\). This turns the
solution set into a cubical complex \(\cplx{\fml}\). For a general complex, the
number of connected components and the Euler characteristic are different
numbers, but for a 2-CNF they are equal:
\[
  \solcomp{\fml}=\euler{\cplx{\fml}} .
\]
The reason is a slicing argument. Setting one variable to \(0\) and to \(1\)
cuts the solution set into two layers, and the assignments that work for both
values of this variable connect the two layers. Because every clause has only
two literals, two different components of this overlap never lie in the same
component of a layer. Hence each component of the overlap joins its own pair of
components, one in each layer.

The second observation is that an Euler characteristic is a local alternating
sum. A subcube is given by a set \(\freeset\subseteq\Vset\) of free variables
together with the values \(\fixasg\in\{0,1\}^{\Vset\setminus\freeset}\) of the
other, fixed variables, and by definition
\[
  \euler{\cplx{\fml}}=\sum_{\freeset\subseteq \Vset}\ \sum_{\fixasg}
  (-1)^{|\freeset|}\,\1\{\text{the subcube given by }\freeset,\fixasg\text{ lies in }\Sol(\fml)\},
\]
that is, we count points, minus segments, plus squares, and so on. This is the
key step: the left-hand side is global, but each indicator on the right is
local, and its expectation is easy to compute. A subcube lies in \(\Sol(\fml)\)
exactly when every clause is satisfied by all of its assignments. For random
signs, the probability that a given clause has this property depends only on
which of its two variables are free, not on the values \(\fixasg\). If we
color the free variables blue and the fixed variables red, then an edge is
satisfied by the whole subcube with probability \(3/4\) when both endpoints are
red, \(1/2\) when exactly one is blue, and \(0\) when both are blue. As a
result, \(\E\,\solcomp{\rform{\Gin}}\) is a signed sum over all red-blue
colorings of the vertices. In each term, every red vertex gives a factor \(2\)
(the two values of a fixed variable), every blue vertex gives a factor \(-1\),
and every edge gives its factor \(3/4\), \(1/2\) or \(0\). Because of cancellation,
this sum can be exponentially small, even though the number of assignments is
exponentially large.

The third observation is that the deletion experiment computes this signed sum.
Delete each edge independently with probability \(1/2\), and give a deleted
edge the weight \(1\) and a kept edge the weight \(1/2\), \(0\) or \(-1\)
according to its color type. The averages are exactly the three numbers
\(3/4\), \(1/2\) and \(0\):
\[
  \tfrac12\bigl(1+\tfrac12\bigr)=\tfrac34,\qquad
  \tfrac12(1+0)=\tfrac12,\qquad
  \tfrac12(1-1)=0 .
\]
We insert this deletion into the sum and exchange it with the sum over
colorings. Once a deletion outcome is fixed, a kept red-blue edge gives \(0\),
so a coloring contributes only if every component of the kept subgraph has a
single color. For one component of the forest, we add the all-red term and the
all-blue term and average over the deferred coins; the result is \(1\) if the
component has no deferred edge, and \(2(3/4)^{t}\) if it has \(t\ge1\) deferred
edges. Taking the product over all components gives
\(2^{\Kdef}(3/4)^{\Tdef}\).

Two features of Theorem~\ref{thm:overview-main-identity} are worth pointing
out. First, it is an equality rather than an inequality, so nothing is lost
when we pass to the right-hand side. Second, the right-hand side does not
involve the formula at all: the question whether \(\rform{\Gin}\) is
satisfiable becomes the question whether, in a random experiment on the graph
\(\Gin\) alone, \(\Kdef\) is small and \(\Tdef\) is large. The only loss in
this part is the inequality \(\solcomp{\fml}\ge1\) for satisfiable \(\fml\).

It remains to prove an upper bound on \(\Kdef\) and a lower bound on \(\Tdef\).

\subsection{The analysis for difficult graphs}\label{subsec:overview-three-conditions}

We now treat difficult graphs, that is, graphs that satisfy \Hyp{} and are
\loctree{\rad}{\nontree}; this is what the reduction above produces. The radius
\(\rad\) and the fraction \(\nontree\) were left open there, and we fix them
here.

Definition~\ref{def:loctree} is used here with two different radii, a large
radius for the bound on \(\Kdef\) and a fixed path length for the bound on
\(\Tdef\); the end of this subsection explains how they fit together.

\begin{theorem}\label{thm:overview-three-conditions}
There are \(\radz=\radz(\degb,\excd)\) and \(\nontreeo=\nontreeo(\degb,\excd)>0\)
such that the following holds. If \(\rad\ge\radz\), \(\nontree\le\nontreeo\), and
the \(\nvert\)-vertex graph \(\Gin\) satisfies \Hyp{} and is
\loctree{\rad}{\nontree}, then
\[
  \E\,\solcomp{\rform{\Gin}}
  =\Ex{2^{\Kdef}\left(\tfrac34\right)^{\Tdef}}
  \le 2e^{-\rate\nvert/2},
\]
where \(\rate=\rate(\degb,\excd)>0\) is the rate of
Lemma~\ref{lem:overview-T-bound}.
\end{theorem}

We prove this in Section~\ref{sec:criterion} from the following two bounds,
which are of very different types: the first is deterministic and purely local,
while the second is probabilistic and comes from the excess.

\begin{lemma}[Upper bound on \(\Kdef\)]\label{lem:overview-K-bound}
If the \(\nvert\)-vertex graph \(\Gin\) is \loctree{\rad}{\nontree}, then
every outcome of the deletion experiment satisfies
\[
  \Kdef\le\left(\nontree+\frac{1}{2\rad}\right)\nvert .
\]
\end{lemma}

A forest component that contains a deferred edge also contains a cycle of
\(\Gin\), namely this edge together with the tree path between its endpoints.
Since different components are disjoint, these cycles are vertex-disjoint, and
hence \(\Kdef\) is at most the largest number of vertex-disjoint cycles in
\(\Gin\). Long cycles, with more than \(2\rad\) vertices, are few: there are
fewer than \(\nvert/(2\rad)\) disjoint ones. A short cycle, of length at most
\(2\rad\), lies inside the radius-\(\rad\) ball of each of its vertices, so a
locally tree-like graph has at most \(\nontree\nvert\) disjoint ones. Thus
\(\Kdef\) is only a small fraction of \(\nvert\) when \(\rad\) is large and
\(\nontree\) is small. Note that this bound uses neither the degree bound nor
the excess.

\begin{lemma}[Exponential moment of \(\Tdef\)]\label{lem:overview-T-bound}
There are an integer \(\plenz=\plenz(\degb,\excd)\),
\(\nontreez=\nontreez(\degb,\excd)>0\) and \(\rate=\rate(\degb,\excd)>0\) such
that the following holds. If \(\Gin\) satisfies \Hyp{} and is
\loctree{\plenz}{\nontreez}, then
\[
  \Ex{\left(\tfrac34\right)^{\Tdef}}\le 2e^{-\rate\nvert} .
\]
\end{lemma}

We state this lemma for \(\Ex{(3/4)^{\Tdef}}\) rather than as a tail bound,
because this is the form used in the proof of
Theorem~\ref{thm:overview-three-conditions}. There the factor \(2^{\Kdef}\) is
at most \(2^{(\nontree+1/(2\rad))\nvert}\), so we only need
\(\Ex{(3/4)^{\Tdef}}\) to decay at a faster exponential rate.

The proof is based on a percolation estimate. For a graph \(H\) and
\(\keepp\in[0,1]\), let \(\perc{H}{\keepp}\) be the random subgraph obtained by
keeping each edge of \(H\) independently with probability \(\keepp\). If we
also flip the coins of the deferred edges, then every edge of \(\Gin\) is kept
independently with probability \(1/2\), and the kept edges form
\(\perc{\Gin}{1/2}\). Since the edges kept in the first stage form a forest,
every cycle of \(\perc{\Gin}{1/2}\) uses a deferred edge, and therefore
\[
  \Tdef\ \ge\ \cyc{\perc{\Gin}{1/2}} .
\]
It is thus enough to show that \(\perc{\Gin}{1/2}\) has linearly many
independent cycles.

This is where the density \(3/2\) appears. Positive excess means
\(\medge\ge(\tfrac32+\excd)\nvert\), so the average degree is at least
\(3+2\excd\). A path that enters a vertex \(v\) can leave it in
\(\dg{\Gin}{v}-1\) ways, and we call this number the branching at \(v\); the
average branching is therefore at least \(2+2\excd\). If each edge is kept with
probability \(\keepp\), long paths with branching \(b\) survive when
\(\keepp\,b>1\) and die out when \(\keepp\,b<1\). The deletion experiment uses
\(\keepp=1/2\), so the condition becomes \(b>2\), which means average degree
above \(3\), or density above \(3/2\). In other words, the threshold in
Theorem~\ref{thm:overview-main} is exactly the density at which percolation
with the fair coin of the deletion experiment becomes supercritical. The
example with three parallel edges fits this picture: it has degree \(3\) and
branching exactly \(2\), which is the critical case.

To get from supercriticality to linearly many cycles, we use one auxiliary
quantity, the probability \(\Lpath{\plen}{H}{\keepp}\) that after percolation
a uniformly random vertex is the start of a path of length \(\plen\). Positive
excess gives a lower bound on this probability by a constant that does not
depend on \(\plen\). A small cycle rank, on the other hand, gives an upper
bound, since a graph with few independent cycles is close to a forest. On a
finite forest of bounded degree, percolation with a fixed retention probability
\(\keepp<1\) rarely keeps a long path from a random vertex. The two bounds can
hold together only if the cycle rank of \(\perc{\Gin}{1/2}\) is linearly large
on average, and a bounded-differences inequality then turns this bound on the
expectation into an exponential bound.

Here the local-tree condition is used differently than in the bound on
\(\Kdef\): it ensures that most vertices have a tree as their neighborhood up
to the distance we need. This makes a second-moment computation on paths
possible, because two paths that separate inside a tree ball never meet again
inside this ball.

\medskip
\noindent\textbf{The order of the constants.} The constants are chosen in the
following order, and each of them depends only on the ones before it.
\begin{enumerate}[label=(\arabic*),leftmargin=3em]
\item \(\excpar\) is given. Lemma~\ref{lem:overview-bounded-core} then fixes
  \(\degb\) and \(\excd\), which depend only on \(\excpar\).
\item Lemma~\ref{lem:overview-T-bound} gives \(\plenz\), \(\nontreez\) and
  \(\rate\), which depend only on \(\degb\) and \(\excd\).
\item Theorem~\ref{thm:overview-three-conditions} chooses \(\radz\) large and
  \(\nontreeo\) small in terms of \(\plenz\), \(\nontreez\) and \(\rate\). These
  are the values of \(\rad\) and \(\nontree\) used in the reduction.
\item Finally, Lemmas~\ref{lem:overview-loop-counting}
  and~\ref{lem:overview-finite-counting}, applied with these values, give
  \(\radp\) and \(\pack\) for Corollary~\ref{cor:overview-packing}.
\end{enumerate}
The important point is step~(2): the rate \(\rate\) does not depend on \(\rad\)
or \(\nontree\). The bound on \(\Kdef\) loses a factor
\(2^{(\nontree+1/(2\rad))\nvert}\), so we need
\((\nontree+\frac1{2\rad})\log2<\rate\). Since \(\rate\) is fixed before
\(\rad\) and \(\nontree\), we can take \(\rad\) large and \(\nontree\) small
enough. We also take \(\radz\ge\plenz\) and \(\nontreeo\le\nontreez\), so that
by Lemma~\ref{lem:loctree-monotone} a graph that is
\loctree{\radz}{\nontreeo} is also \loctree{\plenz}{\nontreez}, and
Lemma~\ref{lem:overview-T-bound} applies.

\subsection{Notation}\label{subsec:notation}

The following table collects the symbols that are used throughout the paper.

%\begin{center}
%\small
%\renewcommand{\arraystretch}{1.15}
%\begin{tabular}{@{}p{0.25\textwidth}p{0.56\textwidth}p{0.13\textwidth}@{}}
%\hline
%Symbol & Meaning & Defined in\\
%\hline
%\(\Gin=(\Vset,\Eset)\), \(\nvert\), \(\medge\)
%  & the input multigraph; \(\nvert=|\Vset|\), \(\medge=|\Eset|\)
%  & \S\ref{subsec:framework}\\
%\(\bvar{v}\), \(\rform{\Gin}\)
%  & the Boolean variable at \(v\); the random 2-CNF formula on \(\Gin\)
%  & \S\ref{subsec:framework}\\
%\(\{\rform{\Gin}\SAT\}\)
%  & the event that \(\rform{\Gin}\) is satisfiable
%  & \S\ref{subsec:framework}\\
%\(\fml\), \(\Sol(\fml)\)
%  & a fixed 2-CNF formula; its set of satisfying assignments, viewed as a
%    graph (the solution graph)
%  & \S\ref{subsec:framework}\\
%\(\ncomp{H}\)
%  & number of components of a graph \(H\); so \(\solcomp{\fml}\) is the number
%    of solution clusters
%  & \S\ref{subsec:framework}\\
%\(\cyc{H}\), \(\exc{H}\)
%  & cycle rank \(|E(H)|-|V(H)|+\ncomp{H}\); excess \(|E(H)|-\frac32|V(H)|\)
%  & \S\ref{subsec:framework}\\
%\(\dg{H}{v}\), \(\dist\)
%  & degree of \(v\) in \(H\); graph distance
%  & \S\ref{subsec:framework}\\
%\(\rad\), \(\ball{\rad}{v}\), \(\ballsz{\rad}\)
%  & a radius; the ball of radius \(\rad\) around \(v\); the bound on its number
%    of vertices
%  & \S\ref{subsec:framework}\\
%\(\degb\)
%  & maximum-degree bound
%  & \S\ref{subsec:framework}\\
%\(\nontree\)
%  & the fraction in ``\loctree{\rad}{\nontree}''
%  & Def.~\ref{def:loctree}\\
%\(\excpar\)
%  & the excess parameter: \(\exc{\Gin}\ge\excpar\nvert\)
%  & Thm.~\ref{thm:main}\\
%\(\excd\), \Hyp
%  & excess density; the standing hypothesis
%  & \S\ref{subsec:overview-graph}\\
%\(\radp\), \(\pack\)
%  & radius of the balls that allow a local contradiction; fraction of vertices
%    with such a ball
%  & Lem.~\ref{lem:overview-loop-counting}, \ref{lem:overview-finite-counting},
%    Cor.~\ref{cor:overview-packing}\\
%\(\freeset\), \(\fixasg\)
%  & the free variables of a subcube; the values of its fixed variables
%  & \S\ref{subsec:overview-main-identity}\\
%\(\cplx{\fml}\), \(\euler{\cdot}\)
%  & the cubical complex of subcubes inside \(\Sol(\fml)\); Euler characteristic
%  & \S\ref{subsec:overview-main-identity}\\
%\(\Ffor\), \(\Tdef\), \(\Kdef\)
%  & the forest of the deletion experiment; number of deferred edges; number of
%    forest components that contain a deferred edge
%  & \S\ref{subsec:overview-main-identity}\\
%\(\perc{H}{\keepp}\), \(\keepp\)
%  & \(H\) with each edge kept independently with probability \(\keepp\);
%    the retention probability
%  & \S\ref{subsec:overview-three-conditions}\\
%\(\plen\), \(\Lpath{\plen}{H}{\keepp}\)
%  & path length; probability that a random vertex starts a kept path of
%    length \(\plen\)
%  & \S\ref{subsec:overview-three-conditions}\\
%\(\radz\), \(\nontreeo\)
%  & the thresholds for \(\rad\) and \(\nontree\)
%  & Thm.~\ref{thm:overview-three-conditions}\\
%\(\plenz\), \(\nontreez\), \(\rate\)
%  & the path length and fraction used for \(\Tdef\); the exponential rate
%  & Lem.~\ref{lem:overview-T-bound}\\
%\(\Prob\), \(\E\)
%  & probability and expectation; a subscript, as in \(\E_{\Phi}\), names the
%    randomness that is averaged over
%  & \\
%\hline
%\end{tabular}
%\end{center}

\subsection{Organization of the paper}

Section~\ref{sec:graph-analysis} contains the graph theory:
Section~\ref{subsec:bounded-core} proves the bounded-degree reduction;
Section~\ref{subsec:unsat-support} proves the unsatisfiable-support lemma, the
positive-excess component reduction and the packing statement,
Corollary~\ref{cor:overview-packing}; and Sections~\ref{subsec:loop-reduction}
and~\ref{subsec:local-tree-reduction} prove the self-loop reduction and the
local-tree reduction. Section~\ref{sec:solution-components} proves the main
identity, Theorem~\ref{thm:overview-main-identity}.
Section~\ref{sec:difficult-graphs} treats the graphs that survive all of the
reductions: it proves the two bounds,
Lemmas~\ref{lem:overview-K-bound} and~\ref{lem:overview-T-bound}, deduces
Theorem~\ref{thm:overview-three-conditions}, and combines everything into
Theorem~\ref{thm:main}. The two estimates on long paths used there are proved
in Appendices~\ref{app:high-functional} and~\ref{app:low-functional}.
